\introduction % Структурный элемент: ВВЕДЕНИЕ

Мир вычислительных технологий за последние десятилетия претерпел
существенные изменения. На смену централизованным и монолитным
приложениям пришли распределённые программные системы, в которых
обработка данных, хранение состояния и выполнение вычислений
осуществляются множеством взаимодействующих узлов. Такой подход
позволяет масштабировать вычислительные ресурсы, повышать доступность
сервисов и обеспечивать устойчивость к частичным отказам оборудования и
сетевой инфраструктуры.

Распределённые системы лежат в основе современных вычислительных
платформ, облачных сервисов, систем хранения данных и транзакционных
приложений. Однако их проектирование и реализация связаны с рядом
фундаментальных трудностей. Поведение такой системы формируется в
условиях асинхронности, конкурентного выполнения, недетерминированной
доставки сообщений и возможных отказов отдельных процессов и каналов
связи. В этих условиях даже небольшая ошибка в алгоритме координации
может приводить к расхождению состояний реплик, потере данных,
некорректному выбору лидера или нарушению свойств безопасности.

Ключевую роль в обеспечении согласованности данных и состояний между
узлами распределённой системы играют алгоритмы консенсуса. Именно они
позволяют множеству процессов прийти к согласованному решению
относительно порядка команд, состояния журнала или результата
взаимодействия, несмотря на задержки сообщений и сбои части узлов.
Среди невизантийских алгоритмов консенсуса наибольшее практическое
значение получили Paxos и Raft. При этом алгоритм Raft представляет
особый интерес благодаря более ясной структуре, декомпозиции на
подзадачи выбора лидера, репликации журнала и обеспечения безопасности,
а также удобству практической реализации.

Для повышения надёжности распределённых алгоритмов всё шире
применяются формальные методы. Одним из наиболее известных и
практически значимых средств формального описания является язык
спецификаций TLA+, позволяющий задавать систему как множество
состояний и допустимых переходов между ними, а также формулировать
инварианты и свойства корректности. Использование средства проверки
моделей TLC делает возможным автоматическое исследование пространства
состояний спецификации и выявление нарушений сформулированных
свойств ещё на этапе проектирования.

Однако наличие корректной формальной спецификации само по себе не
гарантирует корректности программной реализации. Между абстрактной
моделью и исходным кодом существует естественный разрыв. Спецификация
описывает поведение на уровне состояний и переходов, тогда как реальная
система включает сетевое взаимодействие, внутренние буферы,
многопоточность, обработку ошибок, особенности используемых библиотек
и инфраструктурные детали. В процессе развития проекта исходный код
может изменяться быстрее, чем формальная модель, что создаёт риск
расхождения между фактическим поведением программы и поведением,
допускаемым спецификацией.

В связи с этим особую актуальность приобретает задача верификации
соответствия TLA+ спецификации и исходного кода. Практически значимым
подходом к решению данной задачи является валидация трасс исполнения
программы относительно формальной модели. В рамках такого подхода
реализация инструментируется таким образом, чтобы в ходе исполнения
фиксировались наблюдаемые изменения состояния, имеющие смысл в
терминах спецификации. Полученные локальные события объединяются в
общую трассу, после чего проверяется, может ли данная трасса быть
интерпретирована как допустимая последовательность переходов формальной
модели. Тем самым проверяется уже не только корректность самой
спецификации, но и соответствие ей поведения реальной исполняемой
системы.

\subsection*{Обоснование актуальности темы исследования}

Актуальность темы обусловлена одновременно теоретическими и
практическими причинами. С теоретической точки зрения распределённые
системы остаются одной из наиболее сложных областей программной
инженерии, поскольку требуют учёта моделей отказов, ограничений
синхронности, свойств каналов связи и механизмов достижения консенсуса.
С практической точки зрения ошибки в распределённых алгоритмах могут
приводить к серьёзным последствиям: нарушению согласованности данных,
остановке сервисов, потере результатов вычислений и необходимости
дорогостоящей доработки системы на поздних этапах жизненного цикла.

Формальная верификация позволяет существенно повысить уверенность в
корректности проектируемых решений. Вместе с тем наибольшую ценность
она приобретает тогда, когда формальная модель не существует отдельно от
реальной реализации, а используется как критерий проверки фактического
поведения программы. Поэтому разработка методики, обеспечивающей
систематическое сопоставление TLA+ спецификации и исходного кода,
является актуальной задачей, имеющей значение как для исследовательской
деятельности, так и для практики создания надёжного программного
обеспечения.

Дополнительную значимость работе придаёт то, что предложенный подход
рассматривается на примере распределённой системы вычислений,
использующей алгоритм консенсуса Raft. Это позволяет не ограничиваться
абстрактной постановкой задачи, а продемонстрировать применимость
методики в условиях реальной распределённой реализации, включающей
взаимодействие узлов, обмен сообщениями, выбор лидера, репликацию
состояния и обработку отказов.

\subsection*{Цель и задачи работы}

Целью выпускной квалификационной работы является разработка методики
верификации соответствия TLA+ спецификации и исходного кода, а также
её апробация на распределённой системе вычислений, использующей
алгоритм консенсуса Raft.

Для достижения поставленной цели в работе решаются следующие задачи:

\begin{itemize}
    \item проанализировать алгоритмы консенсуса и подходы к верификации
соответствия TLA+ спецификации и исходного кода;
    \item сформулировать и обосновать методику верификации соответствия
спецификации и реализации;
    \item спроектировать и реализовать распределённую систему вычислений,
разработать ее TLA+ спецификацию;
    \item разработать и интегрировать систему верификации;
    \item оценить экономические и организационно-правовые аспекты разработки.
\end{itemize}

\subsection*{Объект и предмет исследования}

Объектом исследования являются распределённые системы вычислений,
использующие алгоритмы консенсуса для обеспечения согласованности
состояния между узлами.

Предметом исследования является методика верификации соответствия
TLA+ спецификации и исходного кода, основанная на инструментировании
реализации, построении трасс исполнения и их последующей проверке
средствами модельной проверки.

\subsection*{Методы исследования}

В работе используются методы теории распределённых систем, формальной
спецификации и модельной проверки, элементы теории переходных систем,
подходы к инструментированию программных реализаций, а также методы
экспериментальной апробации программных средств на тестовых сценариях
работы распределённой системы.

\subsection*{Практическая значимость работы}

Практическая значимость работы заключается в том, что предложенная
методика позволяет обнаруживать расхождения между формальной моделью
и реальным поведением программы на основе анализа трасс исполнения.
Результаты работы могут быть использованы при разработке и
сопровождении распределённых систем, для которых критически важны
согласованность состояния, корректность координации и своевременное
выявление логических ошибок реализации.

Кроме того, разработанные спецификация, программная реализация и
средства интеграции проверки соответствия могут служить основой для
дальнейшего расширения рассматриваемого подхода, включая применение к
другим алгоритмам консенсуса и другим классам распределённых систем.
